% Valores teóricos calculados
\begin{frame}{Magnitude e fase calculadas}
  \small
  \begin{block}{Avaliação de $H(j\omega)$}
    Os valores abaixo foram obtidos com $R_1=R_2=470\,\Omega$,
    $R_3=47\,\mathrm{k}\Omega$ e $C=100\,\mathrm{nF}$.
  \end{block}

  \begin{columns}[T,onlytextwidth]
    \begin{column}{0.48\textwidth}
      \centering
      \begin{tabular}{rcc}
        \toprule
        $f$ (Hz) & $|H(j\omega)|$ & $\phi$ ($^\circ$) \\
        \midrule
        40  & 0,5947  & -90,68 \\
        100 & 1,5431  & -91,77 \\
        200 & 3,5679  & -94,09 \\
        400 & 18,1956 & -111,34 \\
        480 & 49,9733 & 178,13 \\
        \bottomrule
      \end{tabular}
    \end{column}
    \begin{column}{0.48\textwidth}
      \centering
      \begin{tabular}{rcc}
        \toprule
        $f$ (Hz) & $|H(j\omega)|$ & $\phi$ ($^\circ$) \\
        \midrule
        550   & 22,6854 & 116,98 \\
        1100  & 3,7874  & 94,34 \\
        2200  & 1,6149  & 91,85 \\
        5500  & 0,6203  & 90,71 \\
        11000 & 0,3084  & 90,35 \\
        \bottomrule
      \end{tabular}
    \end{column}
  \end{columns}

  \vspace{0.30cm}
  \centering
  \alert{O máximo calculado ocorre próximo de $480\,\mathrm{Hz}$, confirmando a seletividade do filtro.}
\end{frame}